\documentclass[11pt, a4paper]{article}

\usepackage[utf8]{inputenc}
\usepackage[T1]{fontenc}
\usepackage[french]{babel}
\usepackage[margin=2.5cm]{geometry}
\usepackage{booktabs}
\usepackage{longtable}
\usepackage{array}
\usepackage{xcolor}
\usepackage{tikz}
\usepackage{hyperref}
\usepackage{enumitem}
\usepackage{fancyhdr}
\usepackage{titlesec}
\usepackage{parskip}
\usepackage{lmodern}

\usetikzlibrary{shapes.geometric, arrows.meta, positioning}

% Couleurs
\definecolor{bronze}{RGB}{180, 100, 30}
\definecolor{silver}{RGB}{90, 120, 160}
\definecolor{gold}{RGB}{180, 150, 20}
\definecolor{lightgray}{RGB}{245, 245, 245}
\definecolor{darkblue}{RGB}{20, 50, 100}

% En-tête / pied de page
\pagestyle{fancy}
\fancyhf{}
\rhead{\small ENEM -- Volet Ménages T4 2025}
\lhead{\small Pipeline d'apurement}
\cfoot{\thepage}

% Titre de sections
\titleformat{\section}{\large\bfseries\color{darkblue}}{}{0em}{}[\titlerule]
\titleformat{\subsection}{\normalsize\bfseries\color{silver}}{}{0em}{}

% Commande étape
\newcommand{\etape}[2]{%
  \noindent\colorbox{lightgray}{\parbox{\dimexpr\linewidth-2\fboxsep}{%
    \textbf{#1} \hfill \texttt{\small #2}%
  }}%
  \vspace{2pt}%
}

\hypersetup{
  colorlinks=true,
  linkcolor=darkblue,
  urlcolor=silver
}

% ---------------------------------------------------------------
\begin{document}

\begin{titlepage}
  \centering
  \vspace*{2cm}
  {\small\textsc{Institut National de la Statistique -- Côte d'Ivoire}}\\[0.5cm]
  {\small\textsc{Enquête Nationale sur l'Emploi auprès des Ménages}}\\[2cm]
  {\LARGE\bfseries\color{darkblue} Pipeline d'Apurement}\\[0.4cm]
  {\large Volet Ménages -- T4 2025}\\[2cm]
  \rule{0.5\linewidth}{0.4pt}\\[0.8cm]
  {\large Documentation technique}\\[0.4cm]
  {\small Version 1.0 -- Avril 2026}\\[2cm]
  \vfill
  {\small Outil : Stata 17 \quad|\quad Collecte : Survey Solutions \quad|\quad Contrôle de version : Git}
\end{titlepage}

\tableofcontents
\newpage

% ---------------------------------------------------------------
\section{Vue d'ensemble du pipeline}

Le pipeline d'apurement de l'ENEM-Ménages T4 2025 couvre la \textbf{Phase~5 du GSBPM}
(traitement statistique). Il s'insère dans une chaîne de production plus large :

\vspace{0.5cm}
\begin{center}
\begin{tikzpicture}[
  node distance=0.4cm,
  box/.style={rectangle, rounded corners=3pt, draw, minimum width=3.2cm,
              minimum height=0.9cm, align=center, font=\small},
  arr/.style={-{Stealth}, thick}
]
  \node[box, fill=bronze!20, draw=bronze] (col)
    {\textbf{Collecte}\\[-2pt]{\tiny ENE\_SUIVI\_COLLECTE}};
  \node[box, fill=silver!20, draw=silver, right=of col] (apur)
    {\textbf{Apurement}\\[-2pt]{\tiny T4\_2025\_v}};
  \node[box, fill=silver!30, draw=silver, right=of apur] (pond)
    {\textbf{Pondération}\\[-2pt]{\tiny ENE\_SURVEY\_WEIGHTS}};
  \node[box, fill=gold!20, draw=gold, right=of pond] (ind)
    {\textbf{Indicateurs}\\[-2pt]{\tiny ENE\_INDICATORS}};

  \draw[arr, bronze] (col) -- (apur);
  \draw[arr, silver] (apur) -- (pond);
  \draw[arr, gold]   (pond) -- (ind);
\end{tikzpicture}
\end{center}
\vspace{0.3cm}

L'entrée du pipeline est constituée des exports bruts Survey Solutions
(\texttt{datain/brute/}). La sortie est un ensemble de bases standardisées
et apurées par section (\texttt{datain/standard/}), prêtes à être fusionnées
puis transmises au module de pondération.

% ---------------------------------------------------------------
\section{Architecture des fichiers}

\subsection{Dossiers de données}

\begin{longtable}{@{}lp{9cm}@{}}
\toprule
\textbf{Dossier} & \textbf{Contenu} \\
\midrule
\texttt{datain/brute/} & 18 fichiers \texttt{.dta} exports Survey Solutions
  (ménage, membres, revenus, activités, etc.) \\
\texttt{datain/standard/} & Bases découpées par section après
  \texttt{standard\_db.do} \\
\texttt{datatemp/} & Fichiers intermédiaires de codification des emplois \\
\texttt{dataout/} & Bases finales fusionnées (ménage + individus) \\
\texttt{document/} & Fichiers Excel d'erreurs exportés lors de l'apurement \\
\bottomrule
\end{longtable}

\subsection{Fichiers de programme}

\begin{longtable}{@{}lp{8.5cm}@{}}
\toprule
\textbf{Fichier} & \textbf{Rôle} \\
\midrule
\texttt{apurement.do} & Orchestrateur principal -- appelle tous les do-files
  dans l'ordre \\
\texttt{parentfile.do} & Déclaration des chemins globaux (\texttt{\$datain\_brute},
  \texttt{\$dofiles}, etc.) \\
\texttt{Suppression\_entretiens\_T4-2025.do} & Suppression des entretiens invalides
  (138 clés) sur tous les fichiers bruts \\
\texttt{01-traitrement\_autres.do} & Recodage des réponses \og{}autres à préciser\fg{} \\
\texttt{Codification\_T4\_2025.do} & Codification CIAP/CITP des emplois et
  construction de \texttt{branche1}, \texttt{branche2},
  \texttt{code\_citp3\_*} \\
\texttt{Rename\_label\_var.do} & Labellisation des variables binaires
  issues des questions à choix multiples \\
\texttt{standard\_db.do} & Découpage de la base brute en fichiers
  section par section \\
\texttt{01-QM\_S1} à \texttt{S7} & Apurement des sections ménage \\
\texttt{01-INDIV\_S1} à \texttt{S14} & Apurement des sections individu \\
\texttt{Fusion.do / Fusion\_ENE.do} & Fusion des sections en base finale \\
\bottomrule
\end{longtable}

% ---------------------------------------------------------------
\section{Séquence d'exécution}

Le fichier \texttt{apurement.do} exécute les étapes dans l'ordre suivant :

\vspace{0.3cm}
\etape{Étape 1}{Suppression\_entretiens\_T4-2025.do}

Suppression de \textbf{138 entretiens invalides} (vides ou en surnombre)
sur l'ensemble des 18 fichiers \texttt{.dta} de \texttt{datain/brute/}.
Utilise un \texttt{tempfile} de clés et une boucle \texttt{foreach}
avec \texttt{merge~m:1~...~keep(master)}.

\vspace{0.2cm}
\etape{Étape 2}{01-traitrement\_autres.do}

Recodage des champs texte libres (\og{}à préciser\fg{}) en catégories
numériques standardisées (statut, ethnie, religion, assurance, spécialité).

\vspace{0.2cm}
\etape{Étape 3}{Codification\_T4\_2025.do}

\begin{itemize}[topsep=2pt, itemsep=1pt]
  \item Importation des fichiers de codification Excel
  \item Construction de \texttt{code\_citp\_emp\_principal} (6 positions CITP-08)
  \item Construction de \texttt{code\_citp3\_emp\_principal} (3 positions, groupe de base)
  \item Construction de \texttt{branche1} (Agriculture / Industrie / Commerce / Autre service)
  \item Construction de \texttt{branche2} (Primaire / Secondaire / Tertiaire)
\end{itemize}

\vspace{0.2cm}
\etape{Étape 4}{Rename\_label\_var.do}

Labellisation des variables binaires issues des questions à choix multiples
(0 = Non, 1 = Oui).

\vspace{0.2cm}
\etape{Étape 5}{standard\_db.do}

Découpage de la base brute unifiée en fichiers section par section :
\texttt{QM\_S0\_identification}, \texttt{QM\_S7\_logement},
\texttt{INDIV\_S1} à \texttt{INDIV\_S14}, etc.

\vspace{0.2cm}
\etape{Étapes 6--12}{01-QM\_S1 à 01-QM\_S7 \quad(sections ménage)}

Apurement des 7 sections ménage.

\vspace{0.2cm}
\etape{Étapes 13--26}{01-INDIV\_S1 à 01-INDIV\_S14 \quad(sections individu)}

Apurement des 14 sections individu.

% ---------------------------------------------------------------
\section{Méthodologie d'apurement}

Chaque variable est traitée selon une séquence standard en quatre étapes :

\begin{enumerate}[topsep=3pt, itemsep=2pt]
  \item \textbf{Valeurs manquantes} : identification des individus éligibles
    n'ayant pas répondu ; export Excel pour retour enquêteur si nécessaire.
  \item \textbf{Cohérence des sauts} : vérification que seuls les individus
    éligibles ont répondu ; annulation des réponses hors cible.
  \item \textbf{Valeurs permises} : contrôle des plages attendues
    (ex.~: heures travaillées $\leq$ 168, âge cohérent avec niveau d'éducation).
  \item \textbf{Vraisemblance} : détection des incohérences inter-variables
    (ex.~: heures effectives $>$ heures habituelles).
\end{enumerate}

% ---------------------------------------------------------------
\section{Stratégie d'imputation par médiane}

Pour les variables continues à valeurs manquantes résiduelles
(principalement les heures de travail WKT et le loyer L6),
une imputation par \textbf{médiane stratifiée} est appliquée
en trois passes successives, du plus fin au plus grossier :

\vspace{0.3cm}
\begin{center}
\begin{tabular}{clp{7.5cm}}
\toprule
\textbf{Passe} & \textbf{Variables de stratification} & \textbf{Déclencheur} \\
\midrule
1 & \texttt{EP3 branche1 code\_citp3 HH12 HH6}
  & Cellule non vide \\
2 & \texttt{EP3 branche1 HH12 HH6}
  & Toujours manquant après passe 1 \\
3 & \texttt{EP3 branche1}
  & Toujours manquant après passe 2 \\
\bottomrule
\end{tabular}
\end{center}

\vspace{0.3cm}
\begin{description}[topsep=2pt, itemsep=2pt]
  \item[\texttt{EP3}] Statut dans l'emploi (salarié, indépendant, patron, etc.)
  \item[\texttt{branche1}] Branche d'activité économique (Agriculture / Industrie /
    Commerce / Autre service), construite depuis le code CIAP
  \item[\texttt{code\_citp3}] Code CITP-08 à 3 positions (groupe de base),
    dérivé du code 6 positions par \texttt{int(code/1000)}
  \item[\texttt{HH12}] Région administrative
  \item[\texttt{HH6}] Milieu de résidence (1 = Urbain, 2 = Rural)
\end{description}

Le repli progressif garantit qu'aucune valeur manquante ne subsiste
faute de cellule suffisamment peuplée.

% ---------------------------------------------------------------
\section{Gestion de la qualité}

\subsection{Entretiens supprimés}

\begin{tabular}{lrl}
\toprule
\textbf{Motif} & \textbf{Nombre} & \textbf{Critère} \\
\midrule
Entretiens problématiques & 6 & Identifiés manuellement \\
Entretiens vides & 131 & \texttt{HH14 == .} \\
Entretiens en surnombre & 1 & Excédent selon région/ZD \\
\midrule
\textbf{Total} & \textbf{138} & \\
\bottomrule
\end{tabular}

\subsection{Traçabilité des erreurs}

Chaque problème identifié est exporté dans un fichier Excel nommé explicitement
dans \texttt{document/} avant d'être corrigé (ex.~: \texttt{E\_WKT1\_0.xlsx},
\texttt{SRH4\_Manquants.xlsx}, \texttt{Erreur\_L6\_coh.xls}).
Cela permet de justifier chaque décision de correction et de retracer
les anomalies jusqu'à l'enquêteur responsable.

% ---------------------------------------------------------------
\section{Variables construites}

\begin{longtable}{@{}lp{9.5cm}@{}}
\toprule
\textbf{Variable} & \textbf{Description} \\
\midrule
\texttt{EN\_EMP} & Statut d'emploi agrégé (1 = en emploi, 0 = hors emploi) \\
\texttt{branche1} & Branche CIAP agrégée, 4 modalités \\
\texttt{branche2} & Branche CIAP agrégée, 3 secteurs + non classés \\
\texttt{code\_citp3\_emp\_principal} & Code CITP-08 emploi principal, 3 positions \\
\texttt{code\_citp3\_emp\_secondaire} & Code CITP-08 emploi secondaire, 3 positions \\
\texttt{code\_citp3\_pluriactivite1--4} & Codes CITP-08 autres emplois, 3 positions \\
\texttt{NB\_HEURE\_TRAVAIL\_TOTAL} & Heures totales travaillées toutes activités \\
\bottomrule
\end{longtable}

% ---------------------------------------------------------------
\section{Points d'attention pour la prochaine vague}

\begin{enumerate}[topsep=3pt, itemsep=3pt]
  \item \textbf{Interface avec la pondération} : le passage de
    \texttt{T4\_2025\_v} vers \texttt{ENE\_SURVEY\_WEIGHTS} est
    actuellement manuel (copie de fichiers). Une automatisation
    via un script de transfert réduirait le risque d'erreur.
  \item \textbf{Bronze inviolable} : le fichier \texttt{Suppression\_entretiens}
    écrit dans \texttt{datain/brute/}. Pour une architecture médaillon stricte,
    les fichiers bruts devraient être en lecture seule et la suppression
    effectuée à l'entrée de la couche Silver.
  \item \textbf{Fichiers volumineux} : \texttt{01-traitrement\_autres.do}
    (562~Ko) et \texttt{Codification\_T4\_2025.do} (542~Ko) contiennent
    des centaines de \texttt{replace} manuels. Une approche par table de
    correspondance Excel réduirait la maintenance.
  \item \textbf{\texttt{01-INDIV\_S15\_disponibilite.do}} : fichier stub,
    à implémenter ou à retirer définitivement.
\end{enumerate}

\vspace{1cm}
\hrule
\vspace{0.3cm}
{\small\centering
  Document généré depuis l'inspection du code source -- Avril 2026\\
  Pipeline : \texttt{apurement.do} $\rightarrow$ 26 do-files $\rightarrow$
  \texttt{Fusion.do}
\par}

\end{document}
